\pdfoutput=1
\documentclass[11pt]{article}

\usepackage[margin=1in]{geometry}
\usepackage{amsmath,amssymb}
\usepackage{booktabs}
\usepackage{graphicx}
\usepackage{hyperref}
\usepackage{xcolor}
\usepackage{lineno}
\usepackage[numbers,sort&compress]{natbib}

\newcommand{\IIA}{\mathrm{IIA}}
\newcommand{\KL}{D_\mathrm{KL}}
\newcommand{\zc}{z_{\text{causal}}}
\newcommand{\zn}{z_{\text{nuisance}}}
\newcommand{\TODO}[1]{\textcolor{red}{[#1]}}
\renewcommand{\arraystretch}{1.15}

\title{Unconstrained Nonlinear Alignment Trades Reconstruction for Interchange
Accuracy; a Constrained Map Does Not}

\author{Anonymous}
\date{}

\begin{document}
\linenumbers
\maketitle

\begin{abstract}
\citet{sutter2025nonlinear} showed that causal abstraction becomes uninformative
once alignment maps may be arbitrarily nonlinear, reaching 100\% interchange
intervention accuracy (IIA) on randomly initialised language models. We ask what
a \emph{constrained} nonlinear map does under the same tests, which they do not
examine. On six GPT-2 tasks at $k = 1$, linear Distributed Alignment Search
attains strict IIA between $0.00$ and $0.52$ while a label-conditional structured
variational autoencoder attains $\geq 0.90$ on five of six. The two separate more
sharply under three further tests. On a randomly initialised network, the
unconstrained map reaches $\IIA = 0.989$ and the constrained map $0.000$. Under
held-out entities, where the causal variable is token-specific and interchange
should fail, the unconstrained map still reports $1.000$ while the constrained map
reports $0.001$; under held-out templates, where transfer is warranted, it
recovers $0.993$. Measuring round-trip reconstruction explains why: the
unconstrained map loses 50--73\% of the signal decoding its own features, and
penalising that loss reduces its accuracy monotonically across three orders of
magnitude of penalty weight. It obtains accuracy \emph{by means of} unfaithful
reconstruction. The constrained map reconstructs to $0.07\%$ relative error and
attains $\IIA = 1.000$ simultaneously. We also report that our own faithfulness
diagnostic fails: a diversity ratio detects the collapse failure mode on
pretrained models but reads a converged vacuous map as faithful.
\end{abstract}

\section{Introduction}

Causal abstraction asks whether interventions on a model's internal
representation behave like interventions on a high-level algorithmic variable
\citep{geiger2021causal}. Distributed Alignment Search (DAS) operationalises this
by learning an orthogonal rotation identifying a $k$-dimensional subspace,
swapping the in-subspace component between two inputs, and measuring whether the
model produces the counterfactual output \citep{geiger2023finding}.

\citet{sutter2025nonlinear} showed the framework collapses when the alignment map
is unconstrained: sufficiently expressive maps align any network to any
algorithm, and their maps reach 100\% IIA on randomly initialised language
models. That result is both a theorem and an experiment, and the experiment is the
sharper half.

Two facts motivate looking between those extremes. Removing the linearity
constraint makes the framework vacuous. Keeping it leaves linear DAS at strict
IIA $0.00$ on subject--verb agreement in GPT-2, where it shifts probability toward
the counterfactual token without ever making it the argmax. Neither option is
satisfactory, and the question is what constraints short of linearity suffice.

We test a label-conditional structured variational autoencoder: a
label-conditional prior over a supervised causal latent
\citep{khemakhem2020variational}, an unsupervised nuisance latent
\citep{narayanaswamy2017learning}, and an overcomplete causal block, trained with
an interchange term added to the ELBO rather than substituted for it.
\TODO{Confirm this framing is what the final method section says.}

\paragraph{Attribution.} The random-network test is
\citeauthor{sutter2025nonlinear}'s. Measuring IIA on a network with no causal
structure is their contribution, and their result is what establishes that high
accuracy there is uninformative. Our addition is applying it to a constrained map.
The diversity ratio is ours, and its failure is reported as ours.

\section{Results}

\subsection{Six GPT-2 tasks}

\TODO{Table: DAS vs structured VAE, standard and strict IIA at $k=1$, layer 8.
IOI $0.19/0.19$ vs $1.00/1.00$; gender bias $0.62/0.52$ vs $1.00/1.00$;
greater-than $0.93/0.28$ vs $1.00/1.00$; SVA $0.92/0.00$ vs $1.00/1.00$;
hypernymy $0.07/0.02$ vs $0.97/0.90$; capitals $0.12/0.00$ vs $0.51/0.26$.
SINGLE RUN --- needs seeds before this is the headline table.}

The gap between scoring rules is a result in itself. On subject--verb agreement
DAS scores $0.92$ when the counterfactual answer need only outrank the base
answer and $0.00$ when it must be the argmax over the vocabulary. Reporting the
permissive figure alone would describe a method that never produces the
counterfactual token as $92\%$ accurate.

\subsection{Subspace dimension under hard interchange accuracy}

Standard IIA includes pairs where the intervention cannot change the prediction.
Restricting to hard examples (Methods) and sweeping $k$ with five-fold
cross-validation over 1994 IOI pairs:

\begin{table}[t]
\centering\small
\begin{tabular}{lccc}
\toprule
$k$ & DAS & Structured VAE & NL-DAS \\
\midrule
1  & $0.156 \pm 0.360$ & $0.234 \pm 0.424$ & $1.000 \pm 0.000$ \\
2  & $0.273 \pm 0.444$ & $0.696 \pm 0.458$ & $1.000 \pm 0.000$ \\
4  & $0.544 \pm 0.498$ & $\mathbf{0.978 \pm 0.144}$ & $1.000 \pm 0.000$ \\
8  & $0.859 \pm 0.345$ & $0.967 \pm 0.177$ & $1.000 \pm 0.000$ \\
16 & $\mathbf{0.994 \pm 0.074}$ & $0.962 \pm 0.187$ & $1.000 \pm 0.000$ \\
32 & $0.998 \pm 0.030$ & $0.954 \pm 0.205$ & $1.000 \pm 0.000$ \\
\bottomrule
\end{tabular}
\caption{Hard IIA on IOI, mean and standard deviation over five folds. The
structured VAE dominates to $k = 8$; DAS becomes reliable at $k = 16$ and
slightly exceeds it thereafter. NL-DAS reads $1.000$ at every $k$ and every fold,
which \S\ref{sec:random} shows is uninformative.}
\end{table}

The claim is efficiency, not capability: the causal variable is reachable in
about four nonlinear dimensions or about sixteen linear ones. Intervals at low
$k$ are wide because the per-example outcome is near-binary, and the $k=1$ and
$k=2$ comparisons are not separable.

\subsection{Distribution shift}

\TODO{Table: standard / held-out entities / held-out templates on IOI at $k=1$.
NL-DAS $1.000/1.000/0.537$ at $\rho = 0.05$ throughout; structured VAE
$1.000/0.001/0.993$. SINGLE RUN --- needs seeds.}

Held-out entities is a split where failure is the correct outcome, because the
IOI causal variable is token-specific and has no well-defined target across
disjoint entity sets. The unconstrained map transfers perfectly there, for a
damning reason: a map indexed by label never encoded entity identity. Held-out
templates is a split where transfer is warranted, and the ordering reverses.

\subsection{A randomly initialised network}
\label{sec:random}

\TODO{Table: NL-DAS $0.989$, structured VAE $0.000$, five runs across three seeds
and two optimisation budgets. Pre-registered with a withdrawal rule.}

This is the test from \citet{sutter2025nonlinear}. A random network contains no
causal structure, so accuracy above chance is not evidence about the network. The
decision rule was fixed before the run: had the constrained map exceeded $0.30$,
the corresponding claim would have been withdrawn.

\subsection{Reconstruction and interchange accuracy}

Round-trip relative error, $\lVert g(f(h)) - h \rVert^2 / \lVert h \rVert^2$,
measured identically for both maps:

\begin{table}[t]
\centering\small
\begin{tabular}{lccc}
\toprule
& $k{=}2$ & $k{=}8$ & $k{=}32$ \\
\midrule
NL-DAS         & 0.726 & 0.562 & 0.497 \\
Structured VAE & 0.0022 & 0.0009 & 0.0007 \\
\bottomrule
\end{tabular}
\hspace{1em}
\begin{tabular}{lccc}
\toprule
$\lambda_\text{recon}$ & $k{=}2$ & $k{=}8$ & $k{=}32$ \\
\midrule
0     & 0.454 & 0.883 & 0.963 \\
0.01  & 0.343 & 0.874 & 0.972 \\
0.1   & 0.109 & 0.635 & 0.702 \\
1     & 0.128 & 0.630 & 0.698 \\
10    & 0.076 & 0.570 & 0.631 \\
\bottomrule
\end{tabular}
\caption{Left: round-trip relative reconstruction error, three seeds. Right:
NL-DAS interchange accuracy as a function of reconstruction penalty weight.
Penalising reconstruction error reduces accuracy monotonically over three orders
of magnitude.}
\end{table}

The unconstrained map loses between a half and three quarters of the signal
decoding its own features, and forcing it to reconstruct reduces its accuracy at
every penalty weight tested. It therefore obtains accuracy \emph{by means of}
unfaithful reconstruction rather than in spite of it: the interchange succeeds
because the decoder emits activations the model never produces, and closing that
route leaves it with no faithful mechanism. The constrained map reconstructs to
$0.07\%$ relative error and reaches $\IIA = 1.000$ at the same time, so the trade
is not intrinsic to nonlinear alignment.

\subsection{Our faithfulness diagnostic fails}

\TODO{Report: diversity ratio $\rho$ detects collapse on pretrained models
(NL-DAS $\IIA = 1.000$ at $\rho = 0.055$) but reads a converged vacuous map as
faithful ($\rho_\text{within} = 1.068$ at $\IIA = 0.989$ on a random network).
$\rho$ is demoted to a secondary signal. This is why the random-network control is
load-bearing rather than redundant.}

\subsection{Controls}

\TODO{Four ablations against per-task random-subspace floors, which range from
$0.000$ to $0.482$ depending on class count. Every control sits at its floor.}

\section{Discussion}

\TODO{Draft after the tables are final. Must cover: what the constraints buy;
that the failure of our own metric motivates the borrowed control; the
implication that faithfulness and accuracy should be reported together.}

\paragraph{On reporting protocol.}
The DAS paper describes running multiple random initialisations and selecting the
best. A maximum over restarts is an upper order statistic whose expectation grows
with the number of restarts; it is not what a single run yields and is not
comparable across studies using different counts. We report the mean with a
confidence interval as the primary figure and the maximum over seeds alongside it,
for every method, so results can be read under either protocol.
\TODO{Insert the measured inflation: DAS mean $0.919$ against max $0.967$ at
$k=16$ on modular addition.}

\paragraph{Limitations.}
\TODO{Single-run tables above. The mechanism behind $\rho \gg 1$ is unexplained,
including NL-DAS with reconstruction reaching $\rho_\text{within} = 2.249$ on a
pretrained model. The efficiency claim is measured on one model family.}

\section{Methods}

\TODO{Interchange interventions and the two scoring rules. Hard-example
selection. Alignment methods, including that DAS uses random orthogonal
initialisation matching the reference implementation and that the QR
parametrisation used here differs from an explicit orthogonal parametrisation
without affecting the recovered solution. The structured VAE objective. The
diversity ratio. Tasks and controls.}

\section*{Data and Code Availability}
\TODO{Repository and archive.}

\bibliographystyle{plainnat}
\bibliography{references,new_bib_entries}
\end{document}
